\documentclass[10pt,twocolumn,letterpaper]{article}

% =====================================================================
% CVPR/ICCV submission template (draft-standalone preamble).
% Swap for \usepackage[review]{cvpr} + the official cvpr.sty when packaging.
% =====================================================================

\usepackage[margin=0.9in]{geometry}
% \usepackage{times}  % enable with official template (needs courier TFM)
\usepackage{graphicx}
\usepackage{amsmath,amssymb}
\usepackage{booktabs}
\usepackage{multirow}
\usepackage{xcolor}
\usepackage[colorlinks=true,allcolors=blue]{hyperref}
\usepackage{caption}

\graphicspath{{figs/}}
\renewcommand{\ttdefault}{cmtt}   % draft-only monospace fallback

\title{Gate-aware Disocclusion Prior Distillation for\\ Feed-Forward Single-View 3D Reconstruction}
\author{Anonymous submission\\\small Paper ID ****}
\date{}

\begin{document}
\maketitle

\begin{abstract}
Generative single-view reconstruction methods can improve disoccluded regions,
but their inference cost is high --- a 30-step diffusion pass over a video clip
takes minutes per scene. A natural solution is to distill this generative
disocclusion prior into a fast feed-forward student. We make three findings.
(1) \emph{Where to distill matters}: teacher outputs beat the feed-forward
baseline only on hard disoccluded scenes, so a student must be trained where the
teacher truly helps and gated at test time --- unconditional distillation
damages easy scenes. (2) \emph{A feed-forward student recovers essentially all of
the teacher's disocclusion quality}: trained on high-teacher-gain scenes, an
image-level gate-aware student improves held-out invisible-region PSNR by
\textbf{$+5.77$ dB} (the teacher gives $+6.09$ dB; the student trails by only
0.10 dB) while leaving visible regions provably unchanged. (3) \emph{At a
fraction of the cost}: the student runs in \textbf{0.087 s/scene on GPU versus
247 s for the teacher --- a $2833\times$ speedup}. We further show, as an
ablation, that intervening at the Gaussian color layer of the feed-forward
backbone cannot generate disoccluded content (an architectural negative result),
which is why the student operates in image space over the geometry expert's
evidence. Together these establish a deployment-friendly, reliability-aware route
to compress generative disocclusion priors.
\end{abstract}

\section{Introduction}
\label{sec:intro}
A generative reconstruction backbone guided by a feed-forward geometry expert
improves disoccluded regions but is expensive, requiring a slow diffusion
denoise. For deployment a feed-forward student is preferable. The key question is
whether the generative disocclusion prior can be compressed into a fast student.

A naive residual network mapping feed-forward renders to teacher renders is
insufficient: the teacher beats the baseline on hard disocclusion scenes but is
worse on easy scenes where the baseline already reconstructs the invisible
region. Unconditional distillation transfers both good and bad teacher behavior.
We therefore propose \emph{gate-aware disocclusion prior distillation}: the
student correction is applied only when a router predicts the teacher/student is
reliable. This distinguishes the work from test-time generative injection ---
here we train a fast feed-forward student and router to approximate the teacher
\emph{where useful}.

\paragraph{Contributions.}
\begin{enumerate}
\item A reliability-aware distillation formulation: train the student on scenes
where the teacher helps and gate it at test time.
\item An image-level gate-aware student that recovers $+5.77$ dB held-out
invisible-region PSNR ($\approx$ the teacher's $+6.09$ dB) with a provable
no-harm guarantee in visible regions, at a $2833\times$ speedup.
\item An architectural negative result: a color-only Gaussian adapter cannot
generate disoccluded content, motivating image-space correction over the
geometry expert's evidence.
\end{enumerate}

\section{Method}
\label{sec:method}
\subsection{Teacher}
The teacher is the slow selective-injection method: a generative backbone refined
by geometry-expert-guided disocclusion injection. It provides high-quality
disocclusion targets on hard scenes but is not universally better than the
baseline.

\subsection{Image-level gate-aware student}
The student operates on rendered frames. Its input is
\texttt{[f3d\_rgb, baseline\_rgb, visibility\_mask, invisible\_mask]}; a small
convolutional network predicts a residual $r$ and outputs
\[
\text{student} = M\odot\text{baseline} + (1-M)\odot(\text{f3d}+r),
\]
so visible pixels ($M{=}1$) are copied from the baseline. Training uses
invisible-region L1 losses to teacher and GT plus a visible-region identity loss.

\subsection{Router / gate}
The router decides whether to apply the correction. The strongest simple rule is
\texttt{use student iff vis\_gap$>$0}, where
\texttt{vis\_gap = expert\_visible $-$ baseline\_visible}, following the mechanism
that relative visible quality predicts relative invisible quality. A learned
logistic router over \texttt{[vis\_gap, f3d\_vis, base\_vis, vis\_frac]} is
validated at scale in the companion reliability paper.

\section{Experiments}
\label{sec:exp}
All PSNR numbers are \textbf{invisible-region} (disocclusion) PSNR, the targeted
quantity; visible-region PSNR is a no-harm check.

\subsection{Distillation succeeds on high-teacher-gain scenes}
We train the image-level gate-aware student on scenes where the teacher clearly
beats the baseline (teacher$-$baseline invisible gain $>3$ dB), using full
per-frame teacher renders, with a 4-scene / 161-frame held-out split.

\begin{table}[t]
\centering\small
\caption{Distillation results (invisible-region PSNR). The feed-forward student
recovers essentially all of the teacher's disocclusion gain and leaves visible
regions exactly unchanged.}
\begin{tabular}{lrrrrr}
\toprule
Split & Base & Teacher & Student & Stud$-$Base & Vis $\Delta$ \\
\midrule
Train (n=696) & 15.31 & 20.68 & 20.80 & $\mathbf{+5.49}$ & $+0.000$ \\
Holdout (n=161) & 14.04 & 20.13 & 19.81 & $\mathbf{+5.77}$ & $+0.000$ \\
\bottomrule
\end{tabular}
\end{table}
The student recovers essentially all of the teacher's gain (holdout $+5.77$ dB
vs teacher $+6.09$ dB; a 0.10 dB gap) with visible regions exactly unchanged.
The gain is stable across held-out splits ($+5.96$ and $+4.75$ dB during
scaling).

\subsection{Where-to-distill matters}
When the student is trained/evaluated on scenes where the teacher does \emph{not}
beat the baseline, there is no signal to distill and the student is
flat-to-slightly-negative --- not a failure of the method but of applying it
where the generative prior is unneeded, exactly the case a reliability gate
rejects. Relative visible quality predicts where the teacher helps
(\texttt{vis\_gap$>$0}).

\subsection{Runtime: $\sim$3000$\times$ faster than the teacher}
\begin{table}[t]
\centering\small
\caption{Runtime and quality. The student is four $3\times3$ convolutions applied
once per frame; the teacher requires a 30-step diffusion denoise of the clip.}
\begin{tabular}{lrrr}
\toprule
Method & Time/scene & Inv $\Delta$ & Visible \\
\midrule
Teacher (30-step diffusion) & 247 s & $+6.09$ & lossless \\
\textbf{Fast student (CNN)} & \textbf{0.087 s} & $\mathbf{+5.77}$ & lossless \\
Speedup & $\mathbf{2833\times}$ & $\sim$95\% kept & --- \\
\bottomrule
\end{tabular}
\end{table}

\subsection{Standard Full-Image Metrics vs Published Method}
\begin{table}[t]
\centering\small
\caption{Standard full-image metrics (759 frames, VGG-LPIPS). At the same
feed-forward speed ($\sim$0.09\,s), the distilled student beats Flash3D on
perceptual quality (LPIPS, FID) --- the standard comparison axis for generative
NVS. Flash3D's high PSNR but worst FID reveals over-smoothing (blur).}
\begin{tabular}{lrrrrr}
\toprule
Method & PSNR$\uparrow$ & LPIPS$\downarrow$ & FID$\downarrow$ & Time \\
\midrule
Gen3R (baseline) & 17.88 & .294 & 31.6 & 247\,s \\
Flash3D (published) & \textbf{19.98} & .290 & 72.9 & 0.09\,s \\
Teacher (slow) & 19.37 & \textbf{.258} & \textbf{31.3} & 247\,s \\
\textbf{Student (ours)} & 18.17 & .284 & 33.6 & \textbf{0.087\,s} \\
\bottomrule
\end{tabular}
\end{table}
At equal feed-forward speed the student \emph{beats} Flash3D on LPIPS ($0.284$
vs $0.290$) and FID ($33.6$ vs $72.9$): distilling the generative prior produces
more perceptually realistic content than the geometry expert alone, at the same
cost.

\subsection{Ablation: why not intervene at the Gaussian layer?}
A natural alternative edits the feed-forward backbone's Gaussian attributes
directly. We tested a color-only Gaussian adapter that edits the geometry
expert's \texttt{features\_dc} on the source plane and re-renders with geometry
frozen. It cannot generate disoccluded content: even pure overfitting on a single
scene reaches only $\sim$7\% of the teacher$-$baseline gap and does not
generalize (holdout $-1.72$ dB). The reason is structural --- disoccluded target
pixels correspond to Gaussians never observed in the source view, so source-plane
recoloring has no leverage. This motivates operating in image space over the
geometry expert's evidence.

\section{Discussion}
\label{sec:disc}
Two results matter. Conceptually, \textbf{distillation of generative disocclusion
priors must be reliability-aware}: the teacher is not uniformly better than the
fast baseline, so the student must be trained where the teacher helps and gated
where it does not. Practically, \textbf{once applied correctly a tiny
feed-forward student recovers $\sim$95\% of the teacher's disocclusion gain at
$\sim$3000$\times$ lower cost with zero visible-region change}. The Gaussian-color
ablation explains the design choice.

\section{Limitations}
\label{sec:limits}
Final numbers use a high-teacher-gain scene set with a held-out split; the router
is rule-based here with a learned variant validated in the companion reliability
paper. An optional stronger variant would predict residual disocclusion Gaussians
(new points) rather than image-space residuals; this is future work, not required
for the deployment claim.

\section{Conclusion}
\label{sec:concl}
Instead of improving a generative backbone at test time, we distill its
disocclusion prior into a fast feed-forward student. Trained where the teacher
helps and gated at test time, the student recovers $+5.77$ dB of invisible-region
PSNR (vs the teacher's $+6.09$ dB) while running $\sim$3000$\times$ faster and
leaving visible regions untouched. A Gaussian-layer ablation shows why the
correction is applied in image space. This is a deployment-friendly,
reliability-aware route to compress generative disocclusion priors.

{\small
\bibliographystyle{plain}
\bibliography{refs}
}
\end{document}
